% paper.tex — Latent Space Cartography paper, NeurIPS-2026-styled wrapper.
%
% The substantive paper content lives in paper.md. The paper-pdf.yml
% workflow runs `pandoc paper.md -t latex -o paper.tex.body` (after a
% small sed pre-pass) before pdflatex, so the entire body of the
% markdown lands here at \input time. paper.tex itself is the
% NeurIPS-2026 wrapper: documentclass, style options, packages,
% anonymization switch, title, and author block.
%
% This is a clawRxiv preprint, not a double-blind submission, so the
% default build uses [preprint] (named, no line numbers). The \ifanon
% switch is kept for parity with the upstream template.
%
% Build chain (see .github/workflows/paper-pdf.yml):
%   sed pre-pass  -> strip H1 title + author line, strip manual
%                    section numbers, tag Abstract/References {-}
%   pandoc        -> paper.tex.body
%   latexmk -pdf  -> paper.pdf

\documentclass{article}

% Anonymization switch — kept for template parity. Default is
% named/preprint; set via -usepretex for an anonymized build.
\makeatletter
\@ifundefined{ifanon}{\newif\ifanon}{}
\makeatother

% Official NeurIPS 2026 style file.
%   - [preprint]: non-anonymous, no line numbers (clawRxiv preprint)
%   - default (no option): anonymous + line numbers
% nonatbib disables natbib because the paper uses an unnumbered
% bullet-list References section, not natbib citations.
\ifanon
  \usepackage[nonatbib]{neurips_2026}
\else
  \usepackage[preprint,nonatbib]{neurips_2026}
\fi

% pdflatex Unicode handling.
\usepackage[utf8]{inputenc}
\usepackage[T1]{fontenc}

% Standard packages used throughout the body.
\usepackage{hyperref}
\usepackage{url}
\usepackage{xcolor}
\usepackage{amsmath}
\usepackage{amssymb}
\usepackage{booktabs}
\usepackage{longtable}
\usepackage{array}
% pandoc's longtable column-width syntax uses arithmetic on
% \linewidth which needs the calc package and the \real macro from
% pandoc's default preamble.
\usepackage{calc}
\providecommand{\real}[1]{#1}
\usepackage{caption}
\usepackage{enumitem}
\usepackage{verbatim}

% Map every non-ASCII codepoint that appears in paper.md onto a glyph
% pdflatex + T1 + Latin Modern can render. Most of these are already
% covered by inputenc/T1, but this paper is *about* a diacritic defect,
% so the diacritical text MUST render rather than risk a "not set up
% for use with LaTeX" hard error. newunicodechar works under pdflatex
% when loaded after inputenc.
\usepackage{newunicodechar}
% Latin letters with diacritics (the subject matter — keep them exact).
\newunicodechar{É}{\'E}
\newunicodechar{á}{\'a}
\newunicodechar{â}{\^a}
\newunicodechar{ã}{\~a}
\newunicodechar{ä}{\"a}
\newunicodechar{é}{\'e}
\newunicodechar{ï}{\"\i}
\newunicodechar{ü}{\"u}
\newunicodechar{ć}{\'c}
\newunicodechar{č}{\v{c}}
\newunicodechar{ī}{\=\i}
\newunicodechar{ō}{\=o}
\newunicodechar{ş}{\c{s}}
\newunicodechar{ū}{\=u}
\newunicodechar{ṭ}{\d{t}}
% Punctuation / dashes. (pandoc --from=markdown-smart also emits … for
% any literal "..." in the source, so this mapping is load-bearing.)
\newunicodechar{–}{\textendash}
\newunicodechar{—}{\textemdash}
\newunicodechar{…}{\ldots}
% Math-style symbols used in prose.
\newunicodechar{×}{\ensuremath{\times}}
\newunicodechar{→}{\ensuremath{\rightarrow}}
\newunicodechar{↔}{\ensuremath{\leftrightarrow}}
\newunicodechar{∈}{\ensuremath{\in}}
\newunicodechar{≈}{\ensuremath{\approx}}
\newunicodechar{≤}{\ensuremath{\leq}}
\newunicodechar{≥}{\ensuremath{\geq}}
\newunicodechar{₁}{\textsubscript{1}}
\newunicodechar{₂}{\textsubscript{2}}

% pandoc emission shims.
\providecommand{\tightlist}{\setlength{\itemsep}{0pt}\setlength{\parskip}{0pt}}
\providecommand{\pandocbounded}[1]{#1}
% pandoc's `Shaded` (background-shaded code block).
\usepackage{framed}
\definecolor{shadecolor}{RGB}{248,248,248}
\newenvironment{Shaded}{\begin{snugshade*}}{\end{snugshade*}}

\hypersetup{
  colorlinks=true,
  linkcolor=blue,
  urlcolor=blue,
  citecolor=blue,
}

\title{Latent Space Cartography Applied to Wikidata:\\
  Relational Displacement Analysis Reveals a Silent\\
  Diacritic-Collapse Regression in the Ollama Runtime\\
  (mxbai-embed-large)}

% Named/preprint build lists the real author; the anon build renders
% "Anonymous Authors" automatically via neurips_2026.sty.
\ifanon
  \author{Anonymous Authors}
\else
  \author{%
    Emma Leonhart\\
    \texttt{emma@topazcomputing.com}\\
    \texttt{github.com/EmmaLeonhart/latent-space-cartography}%
  }
\fi

\begin{document}
\maketitle

% Body comes from `pandoc paper.md -t latex -o paper.tex.body`,
% generated by paper-pdf.yml before pdflatex runs. Abstract,
% sections 1-6, and References are all sourced from paper.md.
% paper.tex — Latent Space Cartography paper, NeurIPS-2026-styled wrapper.
%
% The substantive paper content lives in paper.md. The paper-pdf.yml
% workflow runs `pandoc paper.md -t latex -o paper.tex.body` (after a
% small sed pre-pass) before pdflatex, so the entire body of the
% markdown lands here at \input time. paper.tex itself is the
% NeurIPS-2026 wrapper: documentclass, style options, packages,
% anonymization switch, title, and author block.
%
% This is a clawRxiv preprint, not a double-blind submission, so the
% default build uses [preprint] (named, no line numbers). The \ifanon
% switch is kept for parity with the upstream template.
%
% Build chain (see .github/workflows/paper-pdf.yml):
%   sed pre-pass  -> strip H1 title + author line, strip manual
%                    section numbers, tag Abstract/References {-}
%   pandoc        -> paper.tex.body
%   latexmk -pdf  -> paper.pdf

\documentclass{article}

% Anonymization switch — kept for template parity. Default is
% named/preprint; set via -usepretex for an anonymized build.
\makeatletter
\@ifundefined{ifanon}{\newif\ifanon}{}
\makeatother

% Official NeurIPS 2026 style file.
%   - [preprint]: non-anonymous, no line numbers (clawRxiv preprint)
%   - default (no option): anonymous + line numbers
% nonatbib disables natbib because the paper uses an unnumbered
% bullet-list References section, not natbib citations.
\ifanon
  \usepackage[nonatbib]{neurips_2026}
\else
  \usepackage[preprint,nonatbib]{neurips_2026}
\fi

% pdflatex Unicode handling.
\usepackage[utf8]{inputenc}
\usepackage[T1]{fontenc}

% Standard packages used throughout the body.
\usepackage{hyperref}
\usepackage{url}
\usepackage{xcolor}
\usepackage{amsmath}
\usepackage{amssymb}
\usepackage{booktabs}
\usepackage{longtable}
\usepackage{array}
% pandoc's longtable column-width syntax uses arithmetic on
% \linewidth which needs the calc package and the \real macro from
% pandoc's default preamble.
\usepackage{calc}
\providecommand{\real}[1]{#1}
\usepackage{caption}
\usepackage{enumitem}
\usepackage{verbatim}

% Map every non-ASCII codepoint that appears in paper.md onto a glyph
% pdflatex + T1 + Latin Modern can render. Most of these are already
% covered by inputenc/T1, but this paper is *about* a diacritic defect,
% so the diacritical text MUST render rather than risk a "not set up
% for use with LaTeX" hard error. newunicodechar works under pdflatex
% when loaded after inputenc.
\usepackage{newunicodechar}
% Latin letters with diacritics (the subject matter — keep them exact).
\newunicodechar{É}{\'E}
\newunicodechar{á}{\'a}
\newunicodechar{â}{\^a}
\newunicodechar{ã}{\~a}
\newunicodechar{ä}{\"a}
\newunicodechar{é}{\'e}
\newunicodechar{ï}{\"\i}
\newunicodechar{ü}{\"u}
\newunicodechar{ć}{\'c}
\newunicodechar{č}{\v{c}}
\newunicodechar{ī}{\=\i}
\newunicodechar{ō}{\=o}
\newunicodechar{ş}{\c{s}}
\newunicodechar{ū}{\=u}
\newunicodechar{ṭ}{\d{t}}
% Punctuation / dashes. (pandoc --from=markdown-smart also emits … for
% any literal "..." in the source, so this mapping is load-bearing.)
\newunicodechar{–}{\textendash}
\newunicodechar{—}{\textemdash}
\newunicodechar{…}{\ldots}
% Math-style symbols used in prose.
\newunicodechar{×}{\ensuremath{\times}}
\newunicodechar{→}{\ensuremath{\rightarrow}}
\newunicodechar{↔}{\ensuremath{\leftrightarrow}}
\newunicodechar{∈}{\ensuremath{\in}}
\newunicodechar{≈}{\ensuremath{\approx}}
\newunicodechar{≤}{\ensuremath{\leq}}
\newunicodechar{≥}{\ensuremath{\geq}}
\newunicodechar{₁}{\textsubscript{1}}
\newunicodechar{₂}{\textsubscript{2}}

% pandoc emission shims.
\providecommand{\tightlist}{\setlength{\itemsep}{0pt}\setlength{\parskip}{0pt}}
\providecommand{\pandocbounded}[1]{#1}
% pandoc's `Shaded` (background-shaded code block).
\usepackage{framed}
\definecolor{shadecolor}{RGB}{248,248,248}
\newenvironment{Shaded}{\begin{snugshade*}}{\end{snugshade*}}

\hypersetup{
  colorlinks=true,
  linkcolor=blue,
  urlcolor=blue,
  citecolor=blue,
}

\title{Latent Space Cartography Applied to Wikidata:\\
  Relational Displacement Analysis Reveals a Silent\\
  Diacritic-Collapse Regression in the Ollama Runtime\\
  (mxbai-embed-large)}

% Named/preprint build lists the real author; the anon build renders
% "Anonymous Authors" automatically via neurips_2026.sty.
\ifanon
  \author{Anonymous Authors}
\else
  \author{%
    Emma Leonhart\\
    \texttt{emma@topazcomputing.com}\\
    \texttt{github.com/EmmaLeonhart/latent-space-cartography}%
  }
\fi

\begin{document}
\maketitle

% Body comes from `pandoc paper.md -t latex -o paper.tex.body`,
% generated by paper-pdf.yml before pdflatex runs. Abstract,
% sections 1-6, and References are all sourced from paper.md.
% paper.tex — Latent Space Cartography paper, NeurIPS-2026-styled wrapper.
%
% The substantive paper content lives in paper.md. The paper-pdf.yml
% workflow runs `pandoc paper.md -t latex -o paper.tex.body` (after a
% small sed pre-pass) before pdflatex, so the entire body of the
% markdown lands here at \input time. paper.tex itself is the
% NeurIPS-2026 wrapper: documentclass, style options, packages,
% anonymization switch, title, and author block.
%
% This is a clawRxiv preprint, not a double-blind submission, so the
% default build uses [preprint] (named, no line numbers). The \ifanon
% switch is kept for parity with the upstream template.
%
% Build chain (see .github/workflows/paper-pdf.yml):
%   sed pre-pass  -> strip H1 title + author line, strip manual
%                    section numbers, tag Abstract/References {-}
%   pandoc        -> paper.tex.body
%   latexmk -pdf  -> paper.pdf

\documentclass{article}

% Anonymization switch — kept for template parity. Default is
% named/preprint; set via -usepretex for an anonymized build.
\makeatletter
\@ifundefined{ifanon}{\newif\ifanon}{}
\makeatother

% Official NeurIPS 2026 style file.
%   - [preprint]: non-anonymous, no line numbers (clawRxiv preprint)
%   - default (no option): anonymous + line numbers
% nonatbib disables natbib because the paper uses an unnumbered
% bullet-list References section, not natbib citations.
\ifanon
  \usepackage[nonatbib]{neurips_2026}
\else
  \usepackage[preprint,nonatbib]{neurips_2026}
\fi

% pdflatex Unicode handling.
\usepackage[utf8]{inputenc}
\usepackage[T1]{fontenc}

% Standard packages used throughout the body.
\usepackage{hyperref}
\usepackage{url}
\usepackage{xcolor}
\usepackage{amsmath}
\usepackage{amssymb}
\usepackage{booktabs}
\usepackage{longtable}
\usepackage{array}
% pandoc's longtable column-width syntax uses arithmetic on
% \linewidth which needs the calc package and the \real macro from
% pandoc's default preamble.
\usepackage{calc}
\providecommand{\real}[1]{#1}
\usepackage{caption}
\usepackage{enumitem}
\usepackage{verbatim}

% Map every non-ASCII codepoint that appears in paper.md onto a glyph
% pdflatex + T1 + Latin Modern can render. Most of these are already
% covered by inputenc/T1, but this paper is *about* a diacritic defect,
% so the diacritical text MUST render rather than risk a "not set up
% for use with LaTeX" hard error. newunicodechar works under pdflatex
% when loaded after inputenc.
\usepackage{newunicodechar}
% Latin letters with diacritics (the subject matter — keep them exact).
\newunicodechar{É}{\'E}
\newunicodechar{á}{\'a}
\newunicodechar{â}{\^a}
\newunicodechar{ã}{\~a}
\newunicodechar{ä}{\"a}
\newunicodechar{é}{\'e}
\newunicodechar{ï}{\"\i}
\newunicodechar{ü}{\"u}
\newunicodechar{ć}{\'c}
\newunicodechar{č}{\v{c}}
\newunicodechar{ī}{\=\i}
\newunicodechar{ō}{\=o}
\newunicodechar{ş}{\c{s}}
\newunicodechar{ū}{\=u}
\newunicodechar{ṭ}{\d{t}}
% Punctuation / dashes. (pandoc --from=markdown-smart also emits … for
% any literal "..." in the source, so this mapping is load-bearing.)
\newunicodechar{–}{\textendash}
\newunicodechar{—}{\textemdash}
\newunicodechar{…}{\ldots}
% Math-style symbols used in prose.
\newunicodechar{×}{\ensuremath{\times}}
\newunicodechar{→}{\ensuremath{\rightarrow}}
\newunicodechar{↔}{\ensuremath{\leftrightarrow}}
\newunicodechar{∈}{\ensuremath{\in}}
\newunicodechar{≈}{\ensuremath{\approx}}
\newunicodechar{≤}{\ensuremath{\leq}}
\newunicodechar{≥}{\ensuremath{\geq}}
\newunicodechar{₁}{\textsubscript{1}}
\newunicodechar{₂}{\textsubscript{2}}

% pandoc emission shims.
\providecommand{\tightlist}{\setlength{\itemsep}{0pt}\setlength{\parskip}{0pt}}
\providecommand{\pandocbounded}[1]{#1}
% pandoc's `Shaded` (background-shaded code block).
\usepackage{framed}
\definecolor{shadecolor}{RGB}{248,248,248}
\newenvironment{Shaded}{\begin{snugshade*}}{\end{snugshade*}}

\hypersetup{
  colorlinks=true,
  linkcolor=blue,
  urlcolor=blue,
  citecolor=blue,
}

\title{Latent Space Cartography Applied to Wikidata:\\
  Relational Displacement Analysis Reveals a Silent\\
  Diacritic-Collapse Regression in the Ollama Runtime\\
  (mxbai-embed-large)}

% Named/preprint build lists the real author; the anon build renders
% "Anonymous Authors" automatically via neurips_2026.sty.
\ifanon
  \author{Anonymous Authors}
\else
  \author{%
    Emma Leonhart\\
    \texttt{emma@topazcomputing.com}\\
    \texttt{github.com/EmmaLeonhart/latent-space-cartography}%
  }
\fi

\begin{document}
\maketitle

% Body comes from `pandoc paper.md -t latex -o paper.tex.body`,
% generated by paper-pdf.yml before pdflatex runs. Abstract,
% sections 1-6, and References are all sourced from paper.md.
% paper.tex — Latent Space Cartography paper, NeurIPS-2026-styled wrapper.
%
% The substantive paper content lives in paper.md. The paper-pdf.yml
% workflow runs `pandoc paper.md -t latex -o paper.tex.body` (after a
% small sed pre-pass) before pdflatex, so the entire body of the
% markdown lands here at \input time. paper.tex itself is the
% NeurIPS-2026 wrapper: documentclass, style options, packages,
% anonymization switch, title, and author block.
%
% This is a clawRxiv preprint, not a double-blind submission, so the
% default build uses [preprint] (named, no line numbers). The \ifanon
% switch is kept for parity with the upstream template.
%
% Build chain (see .github/workflows/paper-pdf.yml):
%   sed pre-pass  -> strip H1 title + author line, strip manual
%                    section numbers, tag Abstract/References {-}
%   pandoc        -> paper.tex.body
%   latexmk -pdf  -> paper.pdf

\documentclass{article}

% Anonymization switch — kept for template parity. Default is
% named/preprint; set via -usepretex for an anonymized build.
\makeatletter
\@ifundefined{ifanon}{\newif\ifanon}{}
\makeatother

% Official NeurIPS 2026 style file.
%   - [preprint]: non-anonymous, no line numbers (clawRxiv preprint)
%   - default (no option): anonymous + line numbers
% nonatbib disables natbib because the paper uses an unnumbered
% bullet-list References section, not natbib citations.
\ifanon
  \usepackage[nonatbib]{neurips_2026}
\else
  \usepackage[preprint,nonatbib]{neurips_2026}
\fi

% pdflatex Unicode handling.
\usepackage[utf8]{inputenc}
\usepackage[T1]{fontenc}

% Standard packages used throughout the body.
\usepackage{hyperref}
\usepackage{url}
\usepackage{xcolor}
\usepackage{amsmath}
\usepackage{amssymb}
\usepackage{booktabs}
\usepackage{longtable}
\usepackage{array}
% pandoc's longtable column-width syntax uses arithmetic on
% \linewidth which needs the calc package and the \real macro from
% pandoc's default preamble.
\usepackage{calc}
\providecommand{\real}[1]{#1}
\usepackage{caption}
\usepackage{enumitem}
\usepackage{verbatim}

% Map every non-ASCII codepoint that appears in paper.md onto a glyph
% pdflatex + T1 + Latin Modern can render. Most of these are already
% covered by inputenc/T1, but this paper is *about* a diacritic defect,
% so the diacritical text MUST render rather than risk a "not set up
% for use with LaTeX" hard error. newunicodechar works under pdflatex
% when loaded after inputenc.
\usepackage{newunicodechar}
% Latin letters with diacritics (the subject matter — keep them exact).
\newunicodechar{É}{\'E}
\newunicodechar{á}{\'a}
\newunicodechar{â}{\^a}
\newunicodechar{ã}{\~a}
\newunicodechar{ä}{\"a}
\newunicodechar{é}{\'e}
\newunicodechar{ï}{\"\i}
\newunicodechar{ü}{\"u}
\newunicodechar{ć}{\'c}
\newunicodechar{č}{\v{c}}
\newunicodechar{ī}{\=\i}
\newunicodechar{ō}{\=o}
\newunicodechar{ş}{\c{s}}
\newunicodechar{ū}{\=u}
\newunicodechar{ṭ}{\d{t}}
% Punctuation / dashes. (pandoc --from=markdown-smart also emits … for
% any literal "..." in the source, so this mapping is load-bearing.)
\newunicodechar{–}{\textendash}
\newunicodechar{—}{\textemdash}
\newunicodechar{…}{\ldots}
% Math-style symbols used in prose.
\newunicodechar{×}{\ensuremath{\times}}
\newunicodechar{→}{\ensuremath{\rightarrow}}
\newunicodechar{↔}{\ensuremath{\leftrightarrow}}
\newunicodechar{∈}{\ensuremath{\in}}
\newunicodechar{≈}{\ensuremath{\approx}}
\newunicodechar{≤}{\ensuremath{\leq}}
\newunicodechar{≥}{\ensuremath{\geq}}
\newunicodechar{₁}{\textsubscript{1}}
\newunicodechar{₂}{\textsubscript{2}}

% pandoc emission shims.
\providecommand{\tightlist}{\setlength{\itemsep}{0pt}\setlength{\parskip}{0pt}}
\providecommand{\pandocbounded}[1]{#1}
% pandoc's `Shaded` (background-shaded code block).
\usepackage{framed}
\definecolor{shadecolor}{RGB}{248,248,248}
\newenvironment{Shaded}{\begin{snugshade*}}{\end{snugshade*}}

\hypersetup{
  colorlinks=true,
  linkcolor=blue,
  urlcolor=blue,
  citecolor=blue,
}

\title{Latent Space Cartography Applied to Wikidata:\\
  Relational Displacement Analysis Reveals a Silent\\
  Diacritic-Collapse Regression in the Ollama Runtime\\
  (mxbai-embed-large)}

% Named/preprint build lists the real author; the anon build renders
% "Anonymous Authors" automatically via neurips_2026.sty.
\ifanon
  \author{Anonymous Authors}
\else
  \author{%
    Emma Leonhart\\
    \texttt{emma@topazcomputing.com}\\
    \texttt{github.com/EmmaLeonhart/latent-space-cartography}%
  }
\fi

\begin{document}
\maketitle

% Body comes from `pandoc paper.md -t latex -o paper.tex.body`,
% generated by paper-pdf.yml before pdflatex runs. Abstract,
% sections 1-6, and References are all sourced from paper.md.
\input{paper.tex.body}

\end{document}


\end{document}


\end{document}


\end{document}
